Pour chacune des affirmations suivantes, indiquer si elle est vraie ou fausse, en justifiant la réponse. Une réponse non justifiée n'est pas prise en compte.

Les deux questions sont indépendantes.

\begin{enumerate}
	\item Dans le plan rapporté à un repère orthonormé, on considère les points
	\[\mathrm{A}(2 ; 0), \quad \mathrm{B}(3 ; 2) \quad \text{et} \quad \mathrm{C}(0 ; 4).\]

	\textbf{Affirmation :} Les droites (AB) et (AC) sont perpendiculaires.

	\item On considère la fonction $f$ définie sur $\R$ par
	\[f(x) = (3x - 1)\e^x.\]

	On note $\mathcal{C}$ sa courbe représentative dans un repère orthogonal du plan.

	On admet que la fonction $f$ est dérivable sur $\R$.

	\begin{enumerate}
		\item \textbf{Affirmation :} la fonction $f$ est croissante sur $\R$.
		\item \textbf{Affirmation :} La tangente à la courbe $\mathcal{C}$ au point d'abscisse $0$ a pour équation
		\[y = 2x - 1.\]
	\end{enumerate}
	\end{enumerate}
